\documentclass[a4paper,10pt]{article}

% =========================================================
% PACKAGES
% =========================================================

% Lab 4 packages
\usepackage{graphicx}
\usepackage{multirow}
\usepackage{array}

% General formatting
\usepackage[margin=1.5cm]{geometry}
\usepackage{longtable}
\usepackage{booktabs}
\usepackage{hyperref}

% If your images are inside a folder called Images,
% you can uncomment the following line:
% \graphicspath{{Images/}}

\title{\textbf{TPT1201 Research Methodology}\\
LaTeX Command Cheat Sheet\\
\large Lab 4}

\author{}
\date{}

\begin{document}

\maketitle
\tableofcontents
\newpage

% =========================================================
% INTRODUCTION
% =========================================================

\section{Introduction}

This document summarizes the important LaTeX commands and examples
covered in Lab 4.

\begin{itemize}
    \item \textbf{Lab 4:} Tables, column alignment, horizontal lines,
    partial horizontal lines, merged columns, merged rows, and images.
\end{itemize}

\textbf{Tip:} In Overleaf, create a new project, open \texttt{main.tex},
delete the existing contents, and paste this entire document.

\newpage

% =========================================================
% LAB 4
% =========================================================

\section{Lab 4 -- Tables and Images}

\subsection{Learning Objectives}

Lab 4 focuses on creating tables and inserting images into a
LaTeX document. :contentReference[oaicite:2]{index=2}

% ---------------------------------------------------------
% LAB 4 CHEAT SHEET
% ---------------------------------------------------------

\subsection{Lab 4 Command Cheat Sheet}

\renewcommand{\arraystretch}{1.4}

\begin{longtable}{p{0.27\textwidth} p{0.32\textwidth} p{0.33\textwidth}}
\caption{Lab 4 -- LaTeX Command Cheat Sheet}
\\
\toprule
\textbf{Command / Code} &
\textbf{Purpose} &
\textbf{Example / Meaning}
\\
\midrule
\endfirsthead

\toprule
\textbf{Command / Code} &
\textbf{Purpose} &
\textbf{Example / Meaning}
\\
\midrule
\endhead

\bottomrule
\endfoot

\texttt{\textbackslash begin\{tabular\}\{...\}} &
Starts a table. &
The column specification is written inside braces.
\\

\texttt{\textbackslash end\{tabular\}} &
Ends the table. &
Marks the end of the tabular environment.
\\

\texttt{c} &
Centers text in a column. &
\texttt{\{c\}} creates one centered column.
\\

\texttt{l} &
Aligns text to the left. &
\texttt{\{l\}} creates one left-aligned column.
\\

\texttt{r} &
Aligns text to the right. &
\texttt{\{r\}} creates one right-aligned column.
\\

\texttt{|} &
Creates a vertical line between columns or at the table boundary. &
\texttt{\{|c|c|\}} creates vertical lines.
\\

\texttt{\&} &
Separates columns in a table row. &
\texttt{Item 1 \& Item 2}.
\\

\texttt{\textbackslash\textbackslash} &
Ends a table row. &
Used after the final cell in a row.
\\

\texttt{\textbackslash hline} &
Creates a horizontal line across the table. &
Useful for separating rows.
\\

\texttt{\textbackslash cline\{i-j\}} &
Creates a partial horizontal line. &
\texttt{\textbackslash cline\{2-3\}} draws a line from column 2 to 3.
\\

\texttt{\textbackslash multicolumn\{n\}\{align\}\{text\}} &
Merges multiple columns into one cell. &
\texttt{\textbackslash multicolumn\{2\}\{|c|\}\{Testing\}}.
\\

\texttt{\textbackslash usepackage\{multirow\}} &
Loads support for merging rows. &
Required before using \texttt{\textbackslash multirow}.
\\

\texttt{\textbackslash multirow\{n\}\{width\}\{text\}} &
Merges multiple rows into one cell. &
\texttt{\textbackslash multirow\{2\}\{*\}\{Already Merged\}}.
\\

\texttt{\textbackslash usepackage\{graphicx\}} &
Loads image insertion commands. &
Required for \texttt{\textbackslash includegraphics}.
\\

\texttt{\textbackslash includegraphics\{file\}} &
Inserts an image. &
\texttt{\textbackslash includegraphics\{mmulogo.png\}}.
\\

\texttt{\textbackslash includegraphics[scale=2.0]\{file\}} &
Inserts an image and scales it. &
Scale 2.0 makes the image twice the default size.
\\

\texttt{\textbackslash graphicspath\{\{Images/\}\}} &
Specifies an image folder. &
Allows images to be stored inside an \texttt{Images} folder.
\\

\texttt{\textbackslash begin\{center\}} &
Centers content. &
Useful for centering an image.
\\

\texttt{\textbackslash end\{center\}} &
Ends the centered section. &
Marks the end of the center environment.
\\

\end{longtable}


% ---------------------------------------------------------
% LAB 4 TABLE EXAMPLES
% ---------------------------------------------------------

\subsection{Lab 4 Table Examples}

\subsubsection{1. One-Column Table}

\begin{verbatim}
\begin{tabular}{|c|}
\hline
row 1 \\
\hline
This is row 2 \\
\hline
Hi, this is another row, row 3 \\
\hline
\end{tabular}
\end{verbatim}

Result:

\begin{center}
\begin{tabular}{|c|}
\hline
row 1 \\
\hline
This is row 2 \\
\hline
Hi, this is another row, row 3 \\
\hline
\end{tabular}
\end{center}

\textbf{Remember:}

\begin{itemize}
    \item \texttt{c} = one centered column.
    \item \texttt{|} = vertical border.
    \item \texttt{\textbackslash hline} = horizontal border.
\end{itemize}

% ---------------------------------------------------------

\subsubsection{2. Two-Column Table}

\begin{verbatim}
\begin{tabular}{|l|l|}
\hline
left item & right side of row 1, name it item 2\\
\hline
row 2 and this is for item 1 & row 2.item2\\
\hline
row 3.item1 & row 3.item2\\
\hline
\end{tabular}
\end{verbatim}

Result:

\begin{center}
\begin{tabular}{|l|l|}
\hline
left item & right side of row 1, name it item 2\\
\hline
row 2 and this is for item 1 & row 2.item2\\
\hline
row 3.item1 & row 3.item2\\
\hline
\end{tabular}
\end{center}

% ---------------------------------------------------------

\subsubsection{3. Three Columns with Different Alignment}

\begin{verbatim}
\begin{tabular}{l c r}
\hline
left aligned & align text in center & make it on right side\\
\hline
row 2 and left for item 1 & center & something\\
\hline
row 3.item1 & center & yet another-right\\
\hline
\end{tabular}
\end{verbatim}

Result:

\begin{center}
\begin{tabular}{l c r}
\hline
left aligned & align text in center & make it on right side\\
\hline
row 2 and left for item 1 & center & something\\
\hline
row 3.item1 & center & yet another-right\\
\hline
\end{tabular}
\end{center}

\textbf{Column alignment:}

\begin{itemize}
    \item \texttt{l} = left
    \item \texttt{c} = center
    \item \texttt{r} = right
\end{itemize}

% ---------------------------------------------------------

\subsubsection{4. Partial Horizontal Lines}

The command \texttt{\textbackslash cline\{i-j\}} creates a horizontal
line only across selected columns.

\begin{verbatim}
\begin{tabular}{|l|c|r|}
\hline
left aligned & align text in center & make it on right side\\
\cline{1-1}
row 2 and left for item 1 & center & something\\
\cline{2-3}
row 3.item1 & center & yet another-right\\
\hline
\end{tabular}
\end{verbatim}

Result:

\begin{center}
\begin{tabular}{|l|c|r|}
\hline
left aligned & align text in center & make it on right side\\
\cline{1-1}
row 2 and left for item 1 & center & something\\
\cline{2-3}
row 3.item1 & center & yet another-right\\
\hline
\end{tabular}
\end{center}

\textbf{Key idea:}

\begin{verbatim}
\cline{1-1}
\end{verbatim}

draws a line only under column 1.

\begin{verbatim}
\cline{2-3}
\end{verbatim}

draws a line under columns 2 through 3.

% ---------------------------------------------------------

\subsubsection{5. Merging Columns with \texttt{\textbackslash multicolumn}}

\begin{verbatim}
\begin{tabular}{|c|c|}
\hline
\multicolumn{2}{|c|}{Testing}\\
\hline
data 1(left column) & data 2 on the second column\\
\hline
\end{tabular}
\end{verbatim}

Result:

\begin{center}
\begin{tabular}{|c|c|}
\hline
\multicolumn{2}{|c|}{Testing}\\
\hline
data 1(left column) & data 2 on the second column\\
\hline
\end{tabular}
\end{center}

\textbf{Syntax:}

\begin{verbatim}
\multicolumn{number of columns}{alignment}{content}
\end{verbatim}

For example:

\begin{verbatim}
\multicolumn{2}{|c|}{Testing}
\end{verbatim}

means that \texttt{Testing} spans across 2 columns and is centered.

% ---------------------------------------------------------

\subsubsection{6. Merging Rows with \texttt{\textbackslash multirow}}

First include:

\begin{verbatim}
\usepackage{multirow}
\end{verbatim}

Then:

\begin{verbatim}
\begin{tabular}{|l|c|c|}
\hline
\multirow{2}{*}{Already Merged} &
\multicolumn{2}{|l|}{Two columns} \\
\cline{2-3}
& 2 & 3 \\
\hline
\end{tabular}
\end{verbatim}

Result:

\begin{center}
\begin{tabular}{|l|c|c|}
\hline
\multirow{2}{*}{Already Merged} &
\multicolumn{2}{|l|}{Two columns} \\
\cline{2-3}
& 2 & 3 \\
\hline
\end{tabular}
\end{center}

\textbf{Syntax:}

\begin{verbatim}
\multirow{number of rows}{width}{content}
\end{verbatim}

% ---------------------------------------------------------

\subsubsection{7. Inserting Images}

First include:

\begin{verbatim}
\usepackage{graphicx}
\end{verbatim}

Then use:

\begin{verbatim}
\begin{center}
\includegraphics[scale=2.0]{mmulogo.png}\\
\end{center}
\end{verbatim}

The image file must be uploaded to your Overleaf project.

For example:

\begin{verbatim}
mmulogo.png
\end{verbatim}

% ---------------------------------------------------------

\subsubsection{8. Images in a Folder}

If your image is stored in a folder named \texttt{Images}, put this
command in the preamble:

\begin{verbatim}
\graphicspath{{Images/}}
\end{verbatim}

Then:

\begin{verbatim}
\documentclass{article}
\usepackage{graphicx}
\graphicspath{{Images/}}

\begin{document}

\begin{center}
\includegraphics[scale=2.0]{mmulogo.png}\\
\end{center}

\end{document}
\end{verbatim}

You do not need to write \texttt{Images/} every time when using
\texttt{\textbackslash includegraphics}.

% =========================================================
% QUICK REVISION
% =========================================================

\newpage

\section{Quick Revision -- What Should I Remember?}

\subsection{Lab 3 Quick Memory Guide}

\begin{center}
\begin{tabular}{p{0.28\textwidth} p{0.62\textwidth}}
\toprule
\textbf{If you want to...} & \textbf{Remember...} \\
\midrule
Write math in a sentence &
Use \texttt{\$...\$} \\

Create a numbered equation &
Use \texttt{equation} \\

Create an unnumbered equation &
Use \texttt{equation*} \\

Make a power &
Use \texttt{\textasciicircum} \\

Make a subscript &
Use \texttt{\_} \\

Make a fraction &
Use \texttt{\textbackslash frac\{top\}\{bottom\}} \\

Make a sum &
Use \texttt{\textbackslash sum} \\

Make a product &
Use \texttt{\textbackslash prod} \\

Create a matrix/array &
Use \texttt{array} \\

Create a piecewise function &
Use \texttt{array} with \texttt{\textbackslash left\{} \\

Create aligned equations &
Use \texttt{IEEEeqnarray} \\

Write an algorithm &
Use \texttt{algorithm} + \texttt{algpseudocode} \\

Create a condition &
Use \texttt{\textbackslash If} \\

Create a loop &
Use \texttt{\textbackslash While} \\

Assign a value &
Use \texttt{\textbackslash gets} \\
\bottomrule
\end{tabular}
\end{center}

\subsection{Lab 4 Quick Memory Guide}

\begin{center}
\begin{tabular}{p{0.28\textwidth} p{0.62\textwidth}}
\toprule
\textbf{If you want to...} & \textbf{Remember...} \\
\midrule
Create a table &
Use \texttt{tabular} \\

Left-align a column &
Use \texttt{l} \\

Center a column &
Use \texttt{c} \\

Right-align a column &
Use \texttt{r} \\

Separate columns &
Use \texttt{\&} \\

Start a new row &
Use \texttt{\textbackslash\textbackslash} \\

Draw a full horizontal line &
Use \texttt{\textbackslash hline} \\

Draw a partial line &
Use \texttt{\textbackslash cline\{i-j\}} \\

Merge columns &
Use \texttt{\textbackslash multicolumn} \\

Merge rows &
Use \texttt{\textbackslash multirow} \\

Insert an image &
Use \texttt{\textbackslash includegraphics} \\

Use images from a folder &
Use \texttt{\textbackslash graphicspath} \\

Center an image &
Use \texttt{center} \\
\bottomrule
\end{tabular}
\end{center}

% =========================================================
% MINI EXAM CHEAT SHEET
% =========================================================

\newpage

\section{Mini Exam Cheat Sheet}

\subsection{Mathematics}

\begin{verbatim}
% Inline math
$a+b=c$

% Numbered equation
\begin{equation}
a^2+b^2=c^2
\end{equation}

% Fraction
\frac{a}{b}

% Power
x^2

% Subscript
x_1

% Sum
\sum_{i=1}^{n}

% Product
\prod_{i=1}^{n}

% Array / Matrix
\begin{array}{ccc}
a & b & c\\
d & e & f
\end{array}
\end{verbatim}

\subsection{Algorithm}

\begin{verbatim}
\usepackage{algorithm}
\usepackage{algpseudocode}

\begin{algorithm}
\caption{MyAlgorithm}
\begin{algorithmic}[1]

\Procedure{MyProcedure}{$x$}

\If{$x>0$}
    \State $x\gets x+1$
\Else
    \State $x\gets 0$
\EndIf

\While{$x<10$}
    \State $x\gets x+1$
\EndWhile

\EndProcedure

\end{algorithmic}
\end{algorithm}
\end{verbatim}

\subsection{Table}

\begin{verbatim}
\begin{tabular}{|c|c|}
\hline
Column 1 & Column 2\\
\hline
Data 1 & Data 2\\
\hline
\end{tabular}
\end{verbatim}

\subsection{Merge Columns}

\begin{verbatim}
\multicolumn{2}{|c|}{Title}
\end{verbatim}

\subsection{Merge Rows}

Remember to add:

\begin{verbatim}
\usepackage{multirow}
\end{verbatim}

Then:

\begin{verbatim}
\multirow{2}{*}{Merged}
\end{verbatim}

\subsection{Image}

Remember to add:

\begin{verbatim}
\usepackage{graphicx}
\end{verbatim}

Then:

\begin{verbatim}
\includegraphics[scale=2.0]{image.png}
\end{verbatim}

% =========================================================

\section{Final Reminder}

\begin{itemize}
    \item \textbf{Lab 4 = Tables + Images}
    \item Always check that required packages are included in the
    preamble.
    \item In tables, \texttt{\&} separates columns and
    \texttt{\textbackslash\textbackslash} creates a new row.
    \item \texttt{\textbackslash hline} creates a full horizontal line,
    while \texttt{\textbackslash cline\{i-j\}} creates a partial line.
    \item \texttt{\textbackslash multicolumn} merges columns.
    \item \texttt{\textbackslash multirow} merges rows.
    \item \texttt{\textbackslash frac} creates fractions.
    \item \texttt{\textasciicircum} creates superscripts and
    \texttt{\_} creates subscripts.
    \item \texttt{algorithm} and \texttt{algpseudocode} are used for
    algorithms.
    \item \texttt{graphicx} is required for inserting images.
\end{itemize}

\end{document}
